\documentclass{ximera}


\ifdefined\HCode
  \newcommand{\alttext}[1]{\HCode{<span class="sr-only">#1</span>}}
\else
  \newcommand{\alttext}[1]{} % Fallback for PDF view
\fi


% MATH -----------------------------------------------------------
\newcommand{\nats}{\mathbb N}
\newcommand{\ints}{\mathbb Z}
\newcommand{\rats}{\mathbb Q}
\newcommand{\reals}{\mathbb R}
\newcommand{\complex}{\mathbb C}
\newcommand{\powerset}{\mathscr P}

%-------------------------------------------------------------

\pgfplotsset{compat=1.18}
\begin{document}
\begin{abstract}
 \end{abstract}

    \title{Class Discussion Activity 25}
    
    \maketitle

\section*{Exercises}

\begin{enumerate}[label=(\arabic*)]

\item Evaluate $\delta(t-1)\ast t^2$ at $t=2$.  

% 1

% Laplace=2e^(-s)/s^3.  Inverse Laplace of that is U(t-1)[(t-1)^2]


\item Suppose $y(t)$ satisfies the integral equation $\displaystyle y(t)+\int_0^t y(x)(x-t)\,\mathrm{d}x=1$.  Determine $y(1)$.  Round to the nearest thousandth.

% 1.543

% Y-Y/s^2=1/s

% s^2Y-Y=s

% (s^2-1)Y=s

% Y=s/(s^2-1)=0.5[1/(s-1)+1/(s+1)]

% y=0.5[e^t+e^{-t}]=cosh(t)

% y=cosh(1)\appprox 1.543




\item Suppose $y(t)$ satisfies the integral equation \\$\displaystyle y(t)+\int_0^t  (t-x)\,y(x)\,dx=\delta{(t-\pi)}$.  Find $y(4)$ to the nearest tenth.


% Y+Y/s^2=e^(- pi s) --> (s^2+1)/s^2 Y=e^(- pi s) ---> Y=e^(- pi s)[(s^2+1)-1]/(s^2+1)=e^(- pi s)[1-1/(s^2+1)] ---> y=U(t-pi){delta(t-pi)+sin(t)}

% sin(4)\approx -0.8




\item Suppose $y(t)$ satisfies the integral equation $\displaystyle y(t)+4\int_0^t  x^2y(t-x)\,\mathrm{d}x=1$.  Determine $y(1)$.  Round to the nearest thousandth.


  Hint:  $s^3+a^3=(s+a)(s^2-as+a^2)$.

    % y(1)=(1/3)[e^(-2)+2e*cos(sqrt(3))]\approx -0.246


    % Y+8Y/s^3=1/s ... Y=s^2/(s^3+8)

    % Y=s^2/(s^3+8)=s^2/[(s+2)(s^2-2s+4)]=(1/3)/(s+2)+((2/3)s-(2/3))/(s^2-2s+4)

    % Y=(1/3)[1/(s+2)+2(s-1)/[(s-1)^2+3]]

    % y=(1/3)[e^(-2t)+2e^t*cos(sqrt(3)t)]



\item Use Euler's method to estimate $y(0.6)$, where $y$ is the unique solution to the IVP $y\,^{\prime}=1- x^y$, $y(0)=1$, using a step size of $\Delta x=0.15$.  Round your estimate to the nearest thousandth. 

% 1.502

    % y(0.15)\approx 1+[1-0^1](0.15)=1.15

    % y(0.3)\approx 1.15+[1-(0.15)^1.15](0.15)\approx 1.28307236808

    % y(0.45)\approx 1.28307236808+[1-(0.3)^(1.28307236808)](0.15)\approx 1.40106868294

    % y(0.6)\approx 1.40106868294+[1-(0.45)^(1.40106868294)](0.15)=1.50206619026




\end{enumerate}

\end{document}